\section{Ballistic exterior-gradient escape and weaker localization rigidity}
\label{sec:r10-dynamics}

We consider the pure mass-critical equation
\begin{equation}\label{eq:r10dyn-equation}
 iu_t=\Delta^2u-|u|^{8/N}u,\qquad N\ge1,
\end{equation}
with the conventions
\[
 p=2+8/N,\quad M=\|u\|_2^2,\quad
 P=\operatorname{Im}\int\overline u\nabla u,\quad
 E=\tfrac12\|\Delta u\|_2^2-\tfrac1p\|u\|_p^p.
\]
Throughout this section, a solution is a standard strong $H^2$ solution
for which these quantities are conserved and the localized virial identity
holds in its integrated energy-space form. Suppose it is forward global and
\begin{equation}\label{eq:r10dyn-bound}
 \sup_{t\ge0}\|u(t)\|_{H^2}\le B<\infty.
\end{equation}
No radial symmetry or finite spatial moment is imposed. Our main result is a
necessary spatial-spreading condition for a hypothetical bounded global
negative-energy solution with zero momentum. It does not establish that
such solutions exist, and it does not exclude all such solutions.

\subsection{A tail estimate and a moving-center identity}

For a center $c(t)$, write
\[
 \begin{gathered}
 D_R^c(t)=\int_{|x-c(t)|>R}|\nabla u(t,x)|^2\,dx,\\
 q_N=\min(4,16/N),\quad
 \kappa_N=\min(1/2,8/N),\quad \beta_N=\min(2,q_N).
 \end{gathered}
\]
The distinction between gradient localization and mass localization is
useful: the former allows noncompact $L^2$ tails whose gradients vanish.

\begin{lemma}[Gradient control of the nonlinear tail]
\label{lem:r10dyn-nonlinear}
For every center $c(t)$ and $R\ge1$,
\begin{equation}\label{eq:r10dyn-nonlineartail}
 \int_{|x-c(t)|>2R}|u(t,x)|^p\,dx
 \le C_{N,M,B}\bigl(D_R^c(t)^{1/2}+R^{-1}\bigr)^{q_N}.
\end{equation}
If $E<0$, there is also $g_*>0$, depending only on $N,M,B,|E|$,
such that $\|\nabla u(t)\|_2\ge g_*$ for all $t\ge0$.
\end{lemma}
\begin{proof}
For $f\in H^2$, homogeneous Sobolev embedding and interpolation give
\begin{equation}\label{eq:r10dyn-interpolation}
 \|f\|_p^p\le
 \begin{cases}
 C_N\|f\|_2^{2(4-N)/N}\|\nabla f\|_2^4,&N\le4,\\[2pt]
 C_N\|\nabla f\|_2^{16/N}\|\Delta f\|_2^{2-8/N},&N>4.
 \end{cases}
\end{equation}
Indeed, the derivative order in the homogeneous embedding into $L^p$ is
$\sigma=N(1/2-1/p)=2N/(N+4)$. Interpolate between orders zero and one
when $\sigma\le1$, and between orders one and two when $\sigma>1$.
This includes the usual $\dot H^1\hookrightarrow L^4$ endpoint at $N=4$.

Choose a smooth cutoff $\eta$ equal to zero on the unit ball and to one
outside the ball of radius two. For
$f(x)=\eta((x-c(t))/R)u(t,x)$, product differentiation gives
\[
 \|f\|_2\le\sqrt M,\quad
 \|\nabla f\|_2\le D_R^c(t)^{1/2}+C\sqrt M/R,\quad
 \|\Delta f\|_2\le C_B.
\]
Substitution in \eqref{eq:r10dyn-interpolation} proves
\eqref{eq:r10dyn-nonlineartail}. Finally, negative energy implies
$\|u(t)\|_p^p=p\|\Delta u(t)\|_2^2/2+p|E|\ge p|E|$.
The global inequalities \eqref{eq:r10dyn-interpolation}, mass conservation,
and \eqref{eq:r10dyn-bound} therefore give the positive gradient lower bound.
\end{proof}

Choose a smooth radial function $\phi_R$ which equals $|x|^2/2$ for
$|x|\le4R$, is constant for $|x|\ge8R$, and satisfies
\[
 D^2\phi_R\le I,\qquad
 \|D^j\phi_R\|_\infty\le C_jR^{2-j}\quad(1\le j\le6).
\]
For a Lipschitz center $a$ with $\|a'\|_\infty\le L$, define
\begin{equation}\label{eq:r10dyn-virialfunctional}
 J_R(t)=2\operatorname{Im}\int
 \nabla\phi_R(x-a(t))\cdot\overline u(t,x)\nabla u(t,x)\,dx.
\end{equation}

\begin{lemma}[Moving virial with gradient-tail error]
\label{lem:r10dyn-virial}
There are constants $C_0=C_0(N,M,B)>0$ and
$C_1=C_1(N,M,B,L)>0$ such that
\begin{align}
 |J_R(t)|&\le C_0R,\label{eq:r10dyn-endpoint}\\
 J_R'(t)&\le16E-2a'(t)\cdot P+
 C_1\bigl(R^{-\beta_N}+D_R^a(t)^{\kappa_N}\bigr).
 \label{eq:r10dyn-virialbound}
\end{align}
The derivative inequality is understood in the integrated sense.
\end{lemma}
\begin{proof}
The endpoint estimate follows directly from Cauchy--Schwarz. Put
$j=\operatorname{Im}(\overline u\nabla u)$ and $\alpha=8/N$.
At a fixed center, differentiation and integration by parts give
\begin{align*}
 J_R'={}&8\operatorname{Re}\int\partial_{jk}\phi_R\,
      \partial_{ij}\overline u\,\partial_{ik}u
 -4\operatorname{Re}\int\partial_{jk}\Delta\phi_R\,
      \partial_j\overline u\,\partial_ku\\
 &-2\int\Delta^2\phi_R|\nabla u|^2
 +\int\Delta^3\phi_R|u|^2
 -\frac{2\alpha}{p}\int\Delta\phi_R|u|^p.
\end{align*}
Repeated indices are summed, and the cutoff is evaluated at $x-a(t)$.
Since $D^2\phi_R\le I$, the leading term is at most
$8\|D^2u\|_2^2=8\|\Delta u\|_2^2$. The higher cutoff derivatives
contribute $O_B(R^{-2})$. The nonlinear term equals
$-16\|u\|_p^p/p$ up to an upper error bounded by
$C\int_{|x-a(t)|>4R}|u|^p$. These principal terms sum to $16E$.

Motion of the center adds exactly
\[
 -2a'(t)\cdot\int D^2\phi_R(x-a(t))j(t,x)\,dx.
\]
Replacing the Hessian by $I$ gives $-2a'\cdot P$. Its error is at most
\[
 C L\int_{|x-a(t)|>4R}|u||\nabla u|
 \le C L\sqrt M\,D_{4R}^a(t)^{1/2}.
\]
Apply Lemma~\ref{lem:r10dyn-nonlinear}, using
$D_{2R}^a,D_{4R}^a\le D_R^a$ and $0\le D_R^a\le B^2$.
The nonlinear-tail exponent is $q_N/2=\min(2,8/N)$; the
center error has exponent $1/2$. Reducing both to
$\kappa_N$ on this bounded interval proves
\eqref{eq:r10dyn-virialbound}. The calculation passes from smooth solutions
to the integrated $H^2$ identity using the displayed bounds. A Lipschitz
center has an almost-everywhere derivative; alternatively, smooth the
center and pass to the integrated identity.
\end{proof}

\subsection{A ballistic necessary alternative}

\begin{theorem}[Persistent exterior gradient at a linear distance]
\label{thm:r10dyn-ballistic}
Suppose \eqref{eq:r10dyn-bound} holds, $E=-e<0$, and $P=0$.
For every Lipschitz center $a:[0,\infty)\to\mathbb R^N$ with
$\|a'\|_\infty\le L$, there are constants $\nu,d>0$, depending only
on $N,M,B,e,L$, such that
\begin{equation}\label{eq:r10dyn-ballistic}
 \liminf_{T\to\infty}\inf_{S\ge0}\frac1T
 \int_S^{S+T}\int_{|x-a(t)|>\nu T}|\nabla u(t,x)|^2\,dx\,dt
 \ge d.
\end{equation}
In particular no choice of bounded-speed center can localize the gradient
inside a sublinearly growing spatial radius in this averaged sense.
\end{theorem}
\begin{proof}
Integrate \eqref{eq:r10dyn-virialbound} on $[S,S+T]$, use $P=0$,
and apply \eqref{eq:r10dyn-endpoint}. Uniformly in $S$,
\begin{equation}\label{eq:r10dyn-integratedescape}
 C_1\frac1T\int_S^{S+T}D_R^a(t)^{\kappa_N}\,dt
 \ge16e-C_1R^{-\beta_N}-2C_0R/T.
\end{equation}
Choose a fixed $\nu>0$ so small that $2C_0\nu\le4e$, and put
$R=\nu T$, held constant during each time integration.
For all sufficiently large $T$, the right side is at least
$8e$. Since $0<\kappa_N\le1$, Jensen's inequality gives
\[
 \left(\frac1T\int_S^{S+T}D_{\nu T}^a(t)\,dt\right)^{\kappa_N}
 \ge\frac1T\int_S^{S+T}D_{\nu T}^a(t)^{\kappa_N}\,dt
 \ge8e/C_1.
\]
This proves \eqref{eq:r10dyn-ballistic} with
$d=(8e/C_1)^{1/\kappa_N}$. The argument also proves the assertion
uniformly over all centers with the specified speed bound.
\end{proof}

\begin{remark}
The theorem asserts persistent exterior gradient in a time-averaged sense,
not at every time. Since $D_R^a(t)\le B^2$, the lower bound also forces a
positive fraction of times in every sufficiently long window to have a
nonzero exterior-gradient amount. The radius $\nu T$ is proportional to
the length of that window. No uniform positive $\nu$ or $d$ is claimed
as $E\to0$ or $B\to\infty$. The exterior gradient is not a physical
dissipation, and the theorem supplies no finite budget for its repeated
transport or replenishment.
\end{remark}

\subsection{Uniform and averaged gradient-localization rigidity}

\begin{theorem}[Averaged localization forces counter-momentum drift]
\label{thm:r10dyn-cesaro}
Suppose \eqref{eq:r10dyn-bound} holds and $E<0$. Let $a$ be Lipschitz,
and suppose
\begin{equation}\label{eq:r10dyn-cesarohyp}
 \lim_{R\to\infty}\limsup_{T\to\infty}\frac1T
 \int_0^TD_R^a(t)\,dt=0.
\end{equation}
Then
\begin{equation}\label{eq:r10dyn-drift}
 \limsup_{T\to\infty}\frac{P\cdot[a(T)-a(0)]}{T}\le8E.
\end{equation}
Consequently \eqref{eq:r10dyn-cesarohyp} is impossible when $P=0$.
\end{theorem}
\begin{proof}
Integrating \eqref{eq:r10dyn-virialbound} and rearranging gives
\[
 2\frac{P\cdot[a(T)-a(0)]}{T}
 \le16E+C_1R^{-\beta_N}+
 C_1\left(\frac1T\int_0^TD_R^a(t)\,dt\right)^{\kappa_N}
 +2C_0R/T.
\]
Here Jensen was again used. First take the upper limit as $T\to\infty$
for fixed $R$, and then let $R\to\infty$. This proves the result.
\end{proof}

\begin{lemma}[Regularizing a gradient concentration center]
\label{lem:r10dyn-center}
Suppose \eqref{eq:r10dyn-bound} holds, $E<0$, and an arbitrary map
$c:[0,\infty)\to\mathbb R^N$ satisfies
\begin{equation}\label{eq:r10dyn-uniformtight}
 \sup_{t\ge0}D_R^c(t)\longrightarrow0\qquad(R\to\infty).
\end{equation}
Then there is a Lipschitz map $a$ with
$\sup_{t\ge0}|a(t)-c(t)|<\infty$. Gradient localization also holds
about $a$.
\end{lemma}
\begin{proof}
The nonlinearity maps bounded $H^2$ sets into bounded $H^{-2}$ sets:
for a test function $h\in H^2$, H\"older and $H^2\hookrightarrow L^p$ give
\[
 \left|\int |u|^{p-2}u\overline h\right|
 \le\|u\|_p^{p-1}\|h\|_p\le C_B\|h\|_{H^2}.
\]
The equation therefore implies a uniform bound on $u_t$ in $H^{-2}$.
Interpolation between $H^{-2}$ and the bounded $H^2$ norm yields
\begin{equation}\label{eq:r10dyn-timeholder}
 \|u(t)-u(s)\|_{H^1}\le C_B|t-s|^{1/4}.
\end{equation}
Let $g_*>0$ be the gradient lower bound from
Lemma~\ref{lem:r10dyn-nonlinear}. Choose $R_0$ so that
$\sup_tD_{R_0}^c(t)\le g_*^2/16$, and choose $\tau>0$ so that
\eqref{eq:r10dyn-timeholder} is at most $g_*/4$ when $|t-s|\le\tau$.

If $|c(t)-c(s)|>2R_0$, the two radius-$R_0$ balls are disjoint.
On $B(c(s),R_0)$, the gradient of $u(s)$ has $L^2$ norm at least
$\sqrt{15}g_*/4$, whereas that of $u(t)$ has norm at most $g_*/4$.
The reverse triangle inequality contradicts
\eqref{eq:r10dyn-timeholder}. Thus
$|c(t)-c(s)|\le2R_0$ whenever $|t-s|\le\tau$.

Let $a$ interpolate linearly through the values $c(n\tau)$,
$n=0,1,\ldots$. Then
\[
 \|a'\|_\infty\le2R_0/\tau,\qquad
 \sup_{t\ge0}|a(t)-c(t)|\le4R_0.
\]
The bounded displacement preserves the uniform tail condition.
\end{proof}

\begin{corollary}[Uniform gradient tightness, with no center regularity]
\label{cor:r10dyn-uniform}
Under \eqref{eq:r10dyn-bound}, $E<0$, and
\eqref{eq:r10dyn-uniformtight},
\[
 \limsup_{T\to\infty}\frac{P\cdot[c(T)-c(0)]}{T}\le8E.
\]
In particular there is no forward-global uniformly $H^2$-bounded
negative-energy zero-momentum solution with uniformly localized gradient
modulo translations. This holds at every mass.
\end{corollary}
\begin{proof}
Apply Lemma~\ref{lem:r10dyn-center}, followed by
Theorem~\ref{thm:r10dyn-cesaro}. The bounded difference between the two
centers does not change the limiting drift quotient.
\end{proof}

Uniform mass tightness and \eqref{eq:r10dyn-bound} imply uniform gradient
tightness by cutoff interpolation. The converse is false for general
bounded $H^2$ families: a fixed localized core can be accompanied by a
fixed-mass packet whose width tends to infinity and whose gradient norm
tends to zero. The result above consequently weakens the required spatial
localization, although it does not weaken the uniform-in-time $H^2$ bound.

\subsection{The unresolved near-critical step}

These results identify an actual requirement on any prospective bounded
global counterexample, but they do not remove the general dynamical gap.
If $M_c$ is the sharp GN critical mass and
$\delta_M=(M/M_c)^{4/N}-1$, the sharp inequality gives
\[
 -\tfrac12\delta_M A(u)\le E(u)<0,\qquad
 A(u)=\|\Delta u\|_2^2.
\]
For fixed negative energy, $A\ge2|E|/\delta_M$. Conversely, normalizing
$A=1$ along a sequence $M\downarrow M_c$ forces $E\to0$, so compactness
can produce an allowed zero-energy optimizer. Normalizing $E=-1$ instead
forces the kinetic norm to diverge. Near-optimizer compactness therefore
does not, by itself, preserve a negative-energy contradiction at the
limiting mass.

A near-critical orbit may have a small, nonzero exterior-gradient budget;
$|E|/A$ can be arbitrarily smaller than that budget. The positive lower
bounds in Theorem~\ref{thm:r10dyn-ballistic} can be compatible with this
small tail. Closing the unrestricted problem requires a further mechanism
that controls sustained radiation or its replenishment, or a construction
showing that such behavior is possible. No such additional mechanism is
asserted here.
